\chapter{Test}

\begin{figure}
\caption{Satz Test}
\fittable{\begin{tabular}{ l l l l |}
    \hline
    \multicolumn{1}{ | l }{\makecell[l]{Ein ehrlicher Mensch ist einer \\ 	
            \small{\textcolor{gray}{Hauptsatz}}}}  &&& \\ \cline{1-1} \sbdashline{2-4}
    
    & \multicolumn{1}{ | l }{\makecell[l]{der etwas verlegen wird \\ 
            \small{\textcolor{gray}{Nebensatz 1}}}}  && \\ \cline{2-2}  \sbdashline{3-4}
    
    && \multicolumn{1}{ | l }{\makecell[l]{wenn man ihm sagt \\ 
            \small{\textcolor{gray}{Nebensatz 2}}}} & \\ \cline{3-3}  \sbdashline{4-4}
    
    &&& \multicolumn{1}{ | l | }{\makecell[l]{er sei ein ehrlicher Mensch \\
             \small{\textcolor{gray}{Nebensatz 3}}}} \\ \cline{4-4}	

\end{tabular}}
\end{figure}

\begin{table}
\begin{tabular}{ll}
\lsptoprule
Test & Tabelle\\\midrule
Wert 1 & Wert 2 \\
Wert 3 & Wert 4\\
\lspbottomrule
\end{tabular}
\caption{Test Tabelle}
\end{table}

\begin{table}[!htbp]
  \centering
  \begin{tabular}{l|l|l}
    \lsptoprule
    {Satz} & \textbf{Hauptsatz} & \textbf{1. Nebensatzebene} \\
    \hline
    \textit{Dass} & \multirow{4}{*}{Subjekt} & -- \\ \cmidrule{3-3}
    \textit{du} & & Subjekt \\ \cmidrule{3-3}
    \textit{mich} & & Akkusativobjekt \\ \cmidrule{3-3}
    \textit{besuchst} & & Prädikat  \\ \cmidrule{2-3}
    \textit{freut} & Prädikat &  \\ \cmidrule{2-2} 
    \textit{mich} & Akkusativobjekt &  \\  
   \lspbottomrule
  \end{tabular}
  \caption{Auflösung des Subjekt-Nebensatzes}
  \label{tab:SubjNS}
\end{table}


\begin{table}[!htbp]
  \centering
  \begin{tabular}{l|l|l}
    %\lsptoprule
    {Satz} & \textbf{Hauptsatz} & \textbf{1. Nebensatzebene} \\
    \hline
    \textit{Dass} & \multirow{4}{*}{Subjekt} & -- \\ \cline{3-3}
    \textit{du} & & Subjekt \\ \cline{3-3}
    \textit{mich} & & Akkusativobjekt \\ \cline{3-3}
    \textit{besuchst} & & Prädikat  \\ \cline{2-3}
    \textit{freut} & Prädikat &  \\ \cline{2-2} 
    \textit{mich} & Akkusativobjekt &  \\  
    %\lspbottomrule
    \hline 
  \end{tabular}
  \caption{Auflösung des Subjekt-Nebensatzes}
  \label{tab:SubjNS}
\end{table}


\begin{table}[!htbp]
  \centering
  \begin{tabular}{l|l|l}
    \lsptoprule
    \textbf{Kategorie} & \textbf{Hilfsverb} & \textbf{Beispiel} \\
    \hline
    Perfekt & \textit{haben} & \textit{Er hat gearbeitet.}\\
    & \textit{sein} & \textit{Sie ist gelaufen.} \\
    Futur & \textit{werden} & \textit{Wir werden tanzen.}\\
    Konjunktiv II & \textit{würde}-Formen & \textit{Wie gerne würde ich tanzen!} \\ 
    Passiv & \textit{werden} & \textit{Der Dieb wird gesucht.}\\  
    & \textit{bekommen}/\textit{kriegen} & \textit{Er bekommt/kriegt das Rad repariert.}\\ 
    \lspbottomrule 
  \end{tabular}
  \caption{Hilfsverben}
  \label{tab:Hilfsverben}
\end{table}



\begin{table}[!htbp]
  \centering
  \begin{tabular}{l|l|l}
    %\lsptoprule
    \hline
    \hline 
    \textbf{Kategorie} & \textbf{Hilfsverb} & \textbf{Beispiel} \\
    \hline
    Perfekt & \textit{haben} & \textit{Er hat gearbeitet.}\\
    & \textit{sein} & \textit{Sie ist gelaufen.} \\
    Futur & \textit{werden} & \textit{Wir werden tanzen.}\\
    Konjunktiv II & \textit{würde}-Formen & \textit{Wie gerne würde ich tanzen!} \\ 
    Passiv & \textit{werden} & \textit{Der Dieb wird gesucht.}\\  
    & \textit{bekommen}/\textit{kriegen} & \textit{Er bekommt/kriegt das Rad repariert.}\\ 
    \hline 
    \hline 
    %\lspbottomrule 
  \end{tabular}
  \caption{Hilfsverben}
  \label{tab:Hilfsverben}
\end{table}